% ============================================================
% CC421 - Inteligencia Artificial
% Tarea-1 PC-1: Fundamentos Éticos en la IA a través de la Filosofía
% Grupo 7 - Kant, Bentham y Mill
% Formato: IEEE doble columna, Times New Roman 10pt
% Referencias: APA 7ma edición
% ============================================================

\documentclass[10pt,twocolumn,letterpaper]{IEEEtran}

% --- Codificación y Fuente ---
\usepackage[utf8]{inputenc}
\usepackage[T1]{fontenc}
\usepackage{times}                  % Times New Roman
\usepackage[spanish]{babel}

% --- Matemáticas y Símbolos ---
\usepackage{amsmath}
\usepackage{amssymb}

% --- Gráficos y Tablas ---
\usepackage{graphicx}
\usepackage{booktabs}
\usepackage{array}
\usepackage{tabularx}

% --- Referencias APA 7 con natbib ---
\usepackage{natbib}
\bibliographystyle{apalike}         % Estilo APA (Author, Year)

% --- Microtipos y Espaciado ---
\usepackage{microtype}
\usepackage{setspace}

% --- Hipervínculos ---
\usepackage[hidelinks]{hyperref}

% --- Configuración de márgenes IEEEtran ---
% IEEEtran gestiona márgenes automáticamente

% ============================================================
% PORTADA / TÍTULO
% ============================================================
\title{\textbf{El Imperativo Categórico frente al Cálculo Utilitarista:\\
Kant, Bentham y Mill en la Arquitectura Ética de la Inteligencia Artificial}}

\author{
  \IEEEauthorblockN{Grupo 7 --- CC421 Inteligencia Artificial}\\
  \IEEEauthorblockA{
    Escuela Profesional de Ciencia de la Computación\\
    Universidad [Nombre de la Universidad]\\
    \textit{\{integrante1, integrante2, integrante3, integrante4, integrante5, integrante6\}@universidad.edu}
  }
}

\begin{document}

\maketitle


% ============================================================
% SECCIONES
% ============================================================

% ============================================================
% SECCIÓN 1: INTRODUCCIÓN
% ============================================================

\section{Introducción}

El 7 de mayo de 2016, un Tesla Model S con el sistema Autopilot activo colisionó fatalmente con un tractor-tráiler en Williston, Florida. La investigación oficial de la NHTSA concluyó que no hubo defecto: el sistema funcionó exactamente como fue diseñado \citep{nhtsa2017pe16007}. El sistema de frenado automático nunca fue concebido para colisiones de trayectoria cruzada ---ese escenario estaba fuera del alcance validado desde 2011. La filosofía moral llevaba siglos debatiendo estos dilemas; la ingeniería de software los había finalmente materializado \citep{mishra2025moral}.

Este accidente no es una anomalía. Es la expresión concreta de una tensión filosófica irresuelta: ¿bajo qué principios debe actuar una máquina cuando todas sus opciones implican daño? La industria ha respondido, mayoritariamente, con cálculos implícitos de utilidad ---priorizando la mayoría de los casos, optimizando para el usuario promedio--- sin reconocer que esta elección constituye ya una postura ética. \citet{mougan2024kantian} evidencian que las métricas de equidad dominantes en IA adoptan supuestos utilitaristas sin declararlo, violando principios de dignidad individual que la ingeniería profesional está obligada a respetar.

El Grupo 7 sostiene la siguiente tesis: \textbf{el imperativo categórico kantiano y el cálculo utilitarista de Bentham y Mill constituyen arquitecturas éticas irreconciliables para el diseño de IA en decisiones críticas}. Su confrontación revela que ningún sistema inteligente puede ser moralmente neutro: toda función objetivo, todo conjunto de reglas de prioridad y todo umbral de decisión incorpora compromisos filosóficos que deben ser explicitados, debatidos y asumidos con responsabilidad profesional. Analizar estos tres pensadores no es un ejercicio académico abstracto; es una precondición para el diseño ético de sistemas que ya toman decisiones de vida o muerte.

El presente artículo se estructura como sigue: la Sección~\ref{sec:conceptual} sintetiza los marcos filosóficos de Kant, Bentham y Mill. La Sección~\ref{sec:aplicacion} los extrapola a la arquitectura de IA y realiza la crítica cruzada entre los tres pensadores. La Sección~\ref{sec:caso} analiza el caso Tesla Autopilot bajo cada lente filosófico. La Sección~\ref{sec:impacto} examina las implicaciones para la responsabilidad profesional del ingeniero. La Sección~\ref{sec:conclusion} sintetiza los hallazgos.


% --- Sección de Desarrollo Conceptual ---
\section{Desarrollo conceptual}

\subsection{Immanuel Kant}

Immanuel Kant nacido en el oriente de prusia en 1724 en una familia de artesanos, encontró una contradicción entre la ciencia clásica, en ese entonces que era mayormente sostenido por la física mecanicista de Issac Newton en el que naturalmente determina en el tiempo y espacio una coherencia consecuente, y la moral la cual Kant sostenía que estaba fundamentalmente construida a partir de la libertad humana, es decir si las reglas del universo están regidas apartir de consecuencias entonces podríamos pensar en que las acciones humanas están estrechamente ligadas a este mundo determinista, por tanto no existiría libertad. Kant estudió meticulosamente esta brecha en su disertación, empezando estudiando la metafísica racionalista de Leibniz y la metafísica racionalista de Christian Wolff, pero los terminó descartando, porque consideró que no toda acción moral debe ser determinada por una doctrina lógica dogmática argumentando que "la moral debe ser motivada en sí misma", es decir que no podemos abstraer las máximas por argumentos sólidos lógicos, porque eso implicaría tratar de perseguir la lógica, dejando de lado la libertad.

\medskip

Al conseguir la disertación, Kant establece el Imperativo Categórico, un método por el cual se puede determinar si una máxima de un acto es moral o inmoral. Establecido por 3 principios:

\medskip

Universalidad: la máxima debe ser aplicada a todos como ley universal.

\medskip

Humanidad: la máxima debe tratar a las personas como fines y nunca considerarlos un medio para otros fines.

\medskip

Autonomía: la máxima debe de considerar a los involucrados como agentes con libre albedrío.

\medskip

Kant concluyó que la ciencia brinda leyes universales y que está sostenido por la percepción racional de lo fenoménico (Realidad consecuente), mientras que lo moral está sostenido por la libertad práctica bajo el imperativo categórico que se encuentra en la realidad nouménica.

\medskip

Para el análisis del caso es necesario considerar los 3 principios de Kant para determinar si el auto autónomo realizó un acto moral o inmoral.

\subsection{Jeremy Bentham: Utilitarismo Clásico y Cálculo Felicífico}
Bentham fundamenta el utilitarismo clásico sobre el principio de utilidad: la acción correcta es aquella que maximiza el bienestar agregado, es decir la mayor felicidad para el mayor número. La premisa es simple, la conducta humana está regida por dos estados: el dolor y el placer \citep{schofield2020jeremy}. 

Para Bentham, la moralidad es una operación aritmética denominada \textit{Cálculo Felicífico}. Este modelo propone que cualquier decisión debe evaluar variables como la intensidad, duración, certeza, proximidad  temporal, fecundidad, pureza y extensión del bienestar generado \citep{quinn2014bentham}. Lo moralmente corecta es la acción que maximiza la suma algebraica de estos factores sobre todos los individuos afectados.

Además, Bentham parte de una concepción hedonista de la naturaleza humana: placer y dolor (como se menciono más arriba), constituyen los criterios últimos de evaluación moral y legislativa. Por ello, su utilitarismo no se limita a juzgar acciones individuales, sino también leyes e instituciones, que deben ser valoradas según su capacidad para aumentar el bienestar general. En este marco, cada persona cuenta por igual dentro del cálculo moral, lo que da al utilitarismo benthamiano un carácter agregativo e imparcial. Sin embargo, esta misma pretensión de cuantificar todos los placeres en una escala común será posteriormente objeto de crítica, especialmente por parte de John Stuart Mill \citep{iepBentham}.

\subsection{John Stuart Mill: Utilitarismo Cualitativo y Justicia}
Mill evoluciona el pensamiento de Bentham al introducir una distinción cualitativa entre los placeres, argumentando que no todos los beneficios tienen el mismo peso ético \citep{west2003introduction}. Como él mismo afirma: ``Es mejor Sócrates insatisfecho que un necio satisfecho'' \citep{mill2014utilitarismo}. Esto implica que un sistema no solo debe buscar la cantidad de bienestar, sino priorizar valores superiores como la justicia y los derechos fundamentales.

Esta distinción tiene una consecuencia directa para el cálculo propuesto por Bentham: si los placeres no son cualitativamente conmensurables, no pueden sumarse en una métrica única. En consecuencia, el cálculo felicífico pierde su fundamento puramente matemático ante la complejidad de los valores cualitativos \citep{mill2014utilitarismo}.

\medskip

Mill complementa esta visión con el \textit{principio del daño} (\textit{harm principle}), principio rector que gobierna la relación entre la sociedad y el individuo. Mill establece que el único fin legítimo para interferir en la libertad de acción de una persona es la autoprotección: ``El único fin para el que el poder puede ejercerse legítimamente sobre un miembro de una comunidad civilizada, en contra de su voluntad, es evitar el daño a otros'' \citep{mill1863utilitarianism}. El bienestar propio del individuo ---físico o moral--- no es razón suficiente para coaccionarlo; sobre sí mismo, sobre su propio cuerpo y mente, el individuo es soberano.

\medskip

Este principio establece una asimetría fundamental: se puede persuadir, razonar o rogar a alguien que cambie su conducta, pero nunca compelerlo a menos que esa conducta cause daño concreto a terceros; solo el daño verificable a otros justifica la intervención \citep{mill2014utilitarismo}.

\medskip

Finalmente, Mill introduce la noción de \textit{principios secundarios}: reglas derivadas de la utilidad que actúan como árbitro cuando obligaciones morales entran en conflicto, evitando recalcular la utilidad en cada situación concreta. Como señala el propio Mill, ``si la utilidad es la fuente última de las obligaciones morales, puede invocarse para decidir entre ellas cuando sus exigencias sean incompatibles'' \citep{mill1863utilitarianism}. Esto constituye el núcleo del \textit{utilitarismo de regla}: en lugar de evaluar acto por acto, el agente ---o sistema--- sigue reglas cuya observancia generalizada maximiza el bienestar colectivo, produciendo un marco normativo estable y predecible.


% TODO
\section{Aplicación de la IA y critica cruzada}
\subsection*{Perspectiva de Kant:}
Para kant un agente autónomo no considera que tenga máxima, esto porque aún no ha sido demostrado que este tenga intenciones, por lo tanto no existiría la máxima, en consecuencia no se podría definir si el agente comete actos morales o inmorales usando el Imperativo Categórico. Entonces la Inteligencia Artificial para Kant no es más que una herramienta en la que estadísticos y matemáticos adjuntan leyes morales en él con intenciones predispuestas en el algoritmo o modelo. 
Para casos de índole ético que interviene la Inteligencia Artificial se debe de juzgar las intenciones y para crear un sistema o arquitectura que usa IA se debe de filtrar las posibles máximas que hacen posible que la IA logre su objetivo, para que con ellos se use los principios del Imperativo Categórico para determinar que la IA, en la toma de decisiones, haga un acto moral. 
\subsection*{}¿Qué pensaría Kant del pensamiento de Mill acerca de un algoritmo de redes sociales que prioriza contenido adictivo porque genera un alto nivel de satisfacción y entretenimiento (placer) en la mayoría de los usuarios?
\subsection*{Al verificar el acto con la intención de mejorar ganancias como fin y usando a la personas como un medio, sería por lo tanto un acto inmoral, debido a que no cumpliría con el segundo principio del Imperativo Categórico que es la Humanidad, las personas no deben ser usadas como un medio sino como un fin solamente. Es por eso que negaría el pensamiento de Mill quien establece el beneficio común.}
\subsection*{}¿Qué pensaría Kant del pensamiento de Bentham acerca de un vehículo autónomo programado para atropellar a un peatón si eso evita un choque donde morirían cuatro pasajeros?
\subsection*{Kant repudiaría la idea de tratar de utilizar a un conjunto de personas como un medio para decidir un acto inmoral, es decir estaría en contra de realizar el algoritmo que determine el acto que cometerá. Sin embargo, sin este algoritmo, Kant estaría de acuerdo en que suceda un evento trágico determinista, la cual análogamente sería otro acto de lo fenoménico.  }
\subsection*{}¿Qué pensaría Kant del pensamiento de Mill acerca de una IA de vigilancia que elimina la privacidad en una ciudad, argumentando que la seguridad resultante produce un bienestar superior para la población?
\subsection*{}
\subsection*{}¿Qué pensaría Kant del pensamiento de Mill acerca de un software de selección de personal que excluye a candidatos brillantes de grupos minoritarios porque su inclusión podría causar fricciones que reduzcan la productividad general de la empresa?
\subsection*{}


% Aplicar el marco teorico a un caso real profundamente documentado.


% ============================================================
% SECCIÓN 5: CRÍTICA Y EVALUACIÓN DE IMPACTO (ALINEACIÓN ABET)
% ============================================================

\section{Crítica y Evaluación de Impacto (Alineación ABET)}
\label{sec:impacto}

El \textit{ABET Student Outcome 4} exige que el ingeniero de computación reconozca responsabilidades profesionales y emita juicios informados basados en principios legales y éticos. El caso Tesla Autopilot revela tres niveles de responsabilidad que los tres marcos filosóficos iluminan desde ángulos distintos.

\textbf{Responsabilidad de diseño (nivel kantiano).} El ingeniero que codifica una función de pérdida para un sistema de VA no está tomando una decisión técnica neutral: está legislando una máxima de acción para millones de interacciones reales. El Código de Ética del IEEE exige que los ingenieros prioricen la seguridad pública sobre cualquier otro objetivo \citep{mishra2025moral}. Desde Kant, esta responsabilidad es categórica: no puede ser cedida a una lógica de mercado ni delegada a la ``objetividad'' del algoritmo. La parálisis deontológica no es una excusa para la inacción; es una señal de que la arquitectura de decisión debe rediseñarse hasta que ninguna de las opciones disponibles viole la dignidad individual.

\textbf{Responsabilidad de transparencia (nivel benthamita).} Bentham exige que el cálculo sea público. \citet{saetra2019tyranny} demuestran que los sistemas de Big Data que operan bajo funciones de optimización opacas generan una ``tiranía de la opinión percibida'': los individuos no pueden evaluar ni contestar los criterios bajo los cuales son clasificados, penalizados u optimizados. Para el VA, esto implica la obligación profesional de publicar los parámetros de la función de utilidad: ¿qué escenarios están cubiertos?, ¿bajo qué condiciones el sistema falla deliberadamente?, ¿quién asume el riesgo residual? La opacidad es, desde el benthamismo, una traición al principio democrático de que cada individuo debe poder calcular su propio riesgo.

\textbf{Responsabilidad de restricción (nivel milleano).} El principio del daño de Mill delimita el espacio de diseño éticamente admisible: el VA puede optimizar para el pasajero, pero no puede externalizar riesgo hacia terceros no consentientes sin su conocimiento. El principio del daño como restricción de diseño tiene una traducción técnica directa: los sistemas de seguridad que protegen a peatones y conductores circundantes no deben ser sujetos de optimización de rendimiento promedio. Son restricciones deontológicas dentro de un marco utilitarista de regla, exactamente como Mill propone \citep{mougan2024kantian}.

El análisis del caso Tesla revela una brecha de accountability que los tres marcos identifican desde ángulos distintos pero convergentes: cuando un sistema de IA mata, la responsabilidad se difumina entre el diseñador del algoritmo, el ingeniero de sistemas, la empresa, el regulador y el usuario. Esta difusión de responsabilidad es filosóficamente inaceptable bajo los tres marcos: Kant exige un agente racional que asuma la máxima del acto; Bentham exige un calculador que asuma la función de utilidad; Mill exige un legislador de reglas que asuma las consecuencias de las normas que impone. En ausencia de accountability clara, los tres sistemas éticos convergen en la misma conclusión: el diseño es moralmente irresponsable.


% Reflexion sintetica de los hallasgos
\section{Conclusiones}


% ============================================================
% REFERENCIAS
% ============================================================
\bibliography{references}

\end{document}
